% [原文 p.39]
% 块首承接：p.38 末式之后接续的那句“……成立．□”是跨页句，由 ch02a 收尾。
特别是，取 $E^0=1$、$p=z$、$E^\infty=L$，则得到
\[
([H]-[1])([L][H]-[1])=0
\]
这是主定理 2.2.1 的一部分。

现在回到线性衔接函数。首先，取有限维向量空间 $E$ 的自同态映射 $T$，再取复平面上一条不经过 $T$ 的任何特征值的圆周 $S$。此时
\[
Q=\frac{1}{2\pi i}\int_S (z-T)^{-1}dz
\]
是 $E$ 的投影算子，并且与 $T$ 交换。因此分解 $E=E_+\oplus E_-$、$E_+=QE$、$E_-=(1-Q)E$ 在 $T$ 之下不变。于是 $T$ 可以写成 $T=T_+\oplus T_-$。这时 $T_+$ 的全部特征值都在 $S$ 的内部，而 $T_-$ 的全部特征值都在 $S$ 的外部。当然，这正是对应于 $S$ 的补集的两个连通分支的 $T$ 的谱分解。

下面把这些结果推广到向量丛。先就记号作一点说明。到此为止，$z$ 与 $p(z)$ 都是 $S$ 上的截面。然而，它们可以自然地扩充为 $L$ 全空间上的截面。在有些论断中，把 $P$ 的无穷远截面一并包含进来会带来方便。因此，例如断言线性衔接 $p(z)=az+b$ 在 $S$ 的外部是同构映射时，也把“$a$ 是同构映射”这一点包含在断言之内。

% [Tutor Note: 易错预警：原书 p.39 定义的两个积分顺序相反，$Q^0$ 用 $p_x^{-1}dp_x$，$Q^\infty$ 用 $dp_xp_x^{-1}$。对非交换的自同态族这两者并不相等，正是积分顺序的差异使 $pQ^0=Q^\infty p$ 成立，抄录时不可把两个积分写成同一个。]

\begin{proposition*}{2.2.8}
  设 $p$ 为相对于 $(E^0,E^\infty)$ 的线性衔接函数，令
  \[
  Q^0_x=\frac{1}{2\pi i}\int_{S_x}p_x^{-1}\,dp_x,\qquad
  Q^\infty_x=\frac{1}{2\pi i}\int_{S_x}dp_x\,p_x^{-1}
  \]
  以此定义 $E^0$ 与 $E^\infty$ 的自同态映射 $Q^0$ 与 $Q^\infty$。此时 $Q^0$ 与 $Q^\infty$ 是投影算子，并且
  \[
  pQ^0=Q^\infty p
  \]
  成立。对 $i=0,\infty$，令 $E^i_+=Q^iE^i$、$E^i_-=(1-Q^i)E^i$，则
% [原文 p.40]
  $E^i=E^i_+\oplus E^i_-$。此时 $p$ 与这些分解相容，即 $p=p_+\oplus p_-$。进一步，$p_+$ 在 $S$ 的外部是同构映射，而 $p_-$ 在 $S$ 的内部是同构映射。
\end{proposition*}

\begin{proof}
  只需对每个点 $x\in X$ 验证论断。换言之，可以假定 $X$ 是单点、$L=\C$，并且 $z$ 只是一个复数。对 $|z|=1$，$p(z)$ 是同构映射，所以可以找到实数 $\alpha>1$，使 $p(\alpha)\colon E^0\to E^\infty$ 为同构映射。为简化计算，用这个同构映射把 $E^0$ 与 $E^\infty$ 等同起来。接着考虑保持单位圆及其内部的共形变换
  \[
  w=\frac{1-\alpha z}{z-\alpha}
  \]
  把它代入 $z$（此处已取 $p(\alpha)=1$），得到
  \[
  p(z)=\frac{w-T}{w+\alpha}
  \]
  其中 $T\in\End(E^0)$。于是
  \begin{align*}
  Q^0&=\frac{1}{2\pi i}\int_{|z|=1}p^{-1}\,dp\\
  &=\frac{1}{2\pi i}\int_{|w|=1}\left(-(w+\alpha)^{-1}dw+(w-T)^{-1}dw\right)\\
  &=\frac{1}{2\pi i}\int_{|w|=1}(w-T)^{-1}dw
  \end{align*}
  成立。最后的等号用到了 $\alpha>1$。同样地，
  \[
  Q^\infty=\frac{1}{2\pi i}\int_{|w|=1}(dw)(w-T)^{-1}=Q^0
  \]
  因此，命题的论断可由关于线性变换 $T$ 的相应论断推出。
\end{proof}

\begin{corollary*}{2.2.9}
  设 $p$ 与命题 2.2.8 中相同，并设
  \[
  p_+=a_+z+b_+,\qquad p_-=a_-z+b_-
  \]
  此时，对
  \[
  p_+(t)=a_+z+tb_+,\qquad p_-(t)=ta_-z+b_-\qquad(0\le t\le1)
  \]
% [原文 p.41]
  令 $p(t)=p_+(t)\oplus p_-(t)$，便得到连接 $p$ 与 $a_+z\oplus b_-$ 的线性衔接函数的同伦。因此
  \[
  (E^0,p,E^\infty)\cong(E^0_+,z,L\otimes E^0_+)\oplus(E^0_-,1,E^0_-)
  \]
\end{corollary*}

\begin{proof}
  由命题 2.2.8 的最后一部分，对 $0\le t\le1$，$p_+(t)$ 与 $p_-(t)$ 在 $S$ 上都是同构映射。因此 $p(t)$ 对 $0\le t\le1$ 都是衔接函数。于是
  \[
  \begin{aligned}
  (E^0,p,E^\infty)&\cong(E^0,p(1),E^\infty)\\
  &\cong(E^0_+,a_+z,E^\infty_+)\oplus(E^0_-,b_-,E^\infty_-)
  \end{aligned}
  \]
  成立。$a_+\colon L\otimes E^0_+\to E^\infty_+$ 与 $b_-\colon E^0_-\to E^\infty_-$ 一定是同构映射，所以
  \[
  (E^0_+,a_+z,E^\infty_+)\cong(E^0_+,z,L\otimes E^0_+)
  \]
  \[
  (E^0_-,b_-,E^\infty_-)\cong(E^0_-,1,E^0_-)
  \]
  由此可知。
\end{proof}

再次考虑次数不超过 $n$ 的多项式衔接函数 $p$。此时令
\[
V^0=\sum_{k=0}^{\infty}L^k\otimes E^0,\qquad
V^\infty=E^\infty\oplus\sum_{k=1}^{n}L^k\otimes E^0
\]
则 $\mathcal{L}^n(p)$ 是相对于 $(V^0,V^\infty)$ 的线性衔接函数。因此，正如前面所述，$\mathcal{L}^n(p)$ 确定分解
\[
V^0=V^0_+\oplus V^0_-
\]
为表明 $V^0_+$ 依赖于 $p$ 与 $n$，记
\[
V^0_+=V_n(E^0,p,E^\infty)
\]
注意这是 $X$ 上的向量丛。设 $p_t$ 是次数不超过 $n$ 的多项式衔接函数的同伦，则在 $X\times I$ 上构造 $V_n$，便得到
% [原文 p.42]
\[
V_n(E^0,p_0,E^\infty)\cong V_n(E^0,p_1,E^\infty)
\]
于是，由证明引理 2.2.6 的 (i) 与 (ii) 时所用的同伦，
\[
V_{n+1}(E^0,p,E^\infty)\cong V_n(E^0,p,E^\infty)
\]
\[
V_{n+1}(L^{-1}\otimes E^0,zp,E^\infty)\cong V_n(E^0,p,E^\infty)\oplus(L^{-1}\otimes E^0)
\]
得到。等价地，也就得到
\[
V_{n+1}(E^0,zp,L\otimes E^\infty)\cong L\otimes V_n(E^0,p,E^\infty)\oplus E^0
\]

把它与推论 2.2.9 以及命题 2.2.5 组合起来，在 $K(P)$ 中得到如下等式。
\[
\begin{aligned}
[E^0,p,E^\infty]+\left\{\sum_{k=1}^{n}[L^k\otimes E^0]\right\}[1]
&=[V_n(E^0,p,E^\infty)][H^{-1}]\\
&\quad+\left\{\sum_{k=0}^{n}[L^k\otimes E^0]-[V_n(E^0,p,E^\infty)]\right\}[1]
\end{aligned}
\]
结果是等式
\[
[E^0,p,E^\infty]=[V_n(E^0,p,E^\infty)]([H^{-1}]-[1])+[E^0][1]
\]
成立。这表明 $[V_n(E^0,p,E^\infty)]\in K(X)$ 完全决定了 $[E^0,p,E^\infty]\in K(P)$。至此可以证明周期性定理了。

% [Tutor Note: 易错预警：上两式里的 $[H^{-1}]$ 就是 $[H]^{-1}$，$H^{-1}$ 是 $H$ 的对偶线丛，$[H][H^{-1}]=[1]$；因此 $[H^{-1}]-[1]$ 才能与前面的 $([H]^{n-1}-[H]^n)$ 对应。原书在 $[H]$ 与 $[H^{-1}]$ 之间来回换写，读时按同一个元素处理。]

设 $t$ 为环 $K(X)$ 上的不定元。此时 $([H]-[1])([L][H]-[1])=0$，所以映射 $t\mapsto[H]$ 诱导出 $K(X)$ 代数同态
\[
\mu\colon K(X)[t]/((t-1)([L]t-1))\longrightarrow K(P)
\]
为证明 $\mu$ 是同构映射，下面具体构造它的逆映射。

先取 $f$ 为相对于 $(E^0,E^\infty)$ 的衔接函数。设 $f_n$ 为其 Fourier 级数的 Cesàro 平均序列，令 $p_n=z^nf_n$。
% [原文 p.43]
此时，当 $n$ 充分大时，$p_n$ 是相对于 $(E^0,L^n\otimes E^\infty)$ 的（次数不超过 $2n$ 的）多项式衔接函数。这里
\[
\nu_n(f)\in K(X)[t]/((t-1)([L]t-1))
\]
由
\[
\nu_n(f)=[V_{2n}(E^0,p_n,L^n\otimes E^\infty)](t^{n-1}-t^n)+[E^0]t^n
\]
定义。

% [Tutor Note: 易错预警：$\nu_n(f)$ 的系数取 $V_{2n}$ 而不是 $V_n$，下标 $2n$ 来自 $p_n=z^nf_n$ 的次数上界 $2n$；后文出现的 $V_{2n+1}$、$V_{2n+2}$ 也据此递推。]

于是，对充分大的 $n$，连接 $p_{n+1}$ 与 $zp_n$ 的线段给出次数不超过 $2(n+1)$ 的多项式衔接函数的同伦。其结果是，由推论 2.2.9 导出的公式得到
\[
\begin{aligned}
V_{2n+2}(E^0,p_{n+1},L^{n+1}\otimes E^\infty)
&\cong V_{2n+2}(E^0,zp_n,L^{n+1}\otimes E^\infty)\\
&\cong V_{2n+1}(E^0,zp_n,L^{n+1}\otimes E^\infty)\\
&\cong L\otimes V_{2n}(E^0,p_n,L^n\otimes E^\infty)\oplus E^0
\end{aligned}
\]
于是，由 $(t-1)([L]t-1)=0$，
\[
\begin{aligned}
\nu_{n+1}(f)&=\{[L][V_{2n}(E^0,p_n,L^n\otimes E^\infty)]+[E^0]\}(t^n-t^{n+1})+[E^0]t^{n+1}\\
&=\nu_n(f)
\end{aligned}
\]

因此，$n$ 充分大时 $\nu_n(f)$ 与 $n$ 无关，只依赖于 $f$。于是把它记作 $\nu(f)$。若 $g$ 与 $f$ 充分接近、$n$ 充分大，则连接 $f_n$ 与 $g_n$ 的线段给出次数不超过 $2n$ 的多项式衔接函数的同伦。所以 $\nu(f)=\nu_n(f)=\nu_n(g)=\nu(g)$。因此 $\nu(f)$ 是 $f$ 的局部常值函数，只依赖于 $f$ 的同伦类。然而，当 $E$ 是 $P$ 上的向量丛、$f$ 是确定 $E$ 的衔接函数时，若定义 $\nu(E)=\nu(f)$，则 $\nu(E)$ 无矛盾地定义，并且只依赖于 $E$ 的同型类。$\nu(E)$ 对直和 $\oplus$ 显然可加，所以诱导出群同态
\[
\nu\colon K(P)\longrightarrow K(X)[t]/((t-1)([L]t-1))
\]
由构造可知，这显然是 $K(X)$ 模同态。

% [原文 p.44]
先确认 $\mu\circ\nu$ 是恒等映射。沿用此前的记号，
\[
\begin{aligned}
\mu\circ\nu(E)&=\mu\{[V_{2n}(E^0,p_n,L^n\otimes E^\infty)](t^{n-1}-t^n)+[E^0]t^n\}\\
&=[V_{2n}(E^0,p_n,L^n\otimes E^\infty)]([H]^{n-1}-[H]^n)+[E^0][H]^n\\
&=[E^0,p_n,L^n\otimes E^\infty][H]^n\\
&=[E^0,f_n,E^\infty]\\
&=[E^0,f,E^\infty]\\
&=[E]
\end{aligned}
\]
于是得到 $\mu\circ\nu(E)=[E]$。由于 $K(P)$ 由形如 $[E]$ 的元素加法生成，$\mu\circ\nu$ 是恒等映射一事便得到证明。

最后证明 $\nu\circ\mu$ 是恒等映射。$\nu\circ\mu$ 是 $K(X)$ 模同态，所以只需对一切 $n\ge0$ 证明 $\nu\circ\mu(t^n)=t^n$。而
\[
\begin{aligned}
\nu\circ\mu(t^n)&=\nu(H^n)\\
&=\nu[1,z^{-n},L^{-n}]\\
&=[V_{2n}(1,1,1)](t^{n-1}-t^n)+[1]t^n\\
&=t^n
\end{aligned}
\]
最后等号用到了 $V_{2n}(1,1,1)=0$。

\section{$K_G(X)$}
\takeaway{把 $K$ 群搬到有限群 $G$ 作用的情形：对 $G$-向量丛的同型类半群 $\Vect_G(X)$ 作 Grothendieck 构造得 $K_G(X)$，$G$ 平凡时回到 $K(X)$，$X$ 为单点时化为特征标环 $R(G)$；§2.2 的全部论证对 $G$ 等变版本同样有效，周期性定理因此对 $K_G$ 成立。}

设 $G$ 为有限群，$X$ 为 $G$ 空间。用 $\Vect_G(X)$ 表示 $X$ 上 $G$-向量丛的同型类的集合。它在直和 $\oplus$ 之下是交换半群。再作 $\Vect_G(X)$ 所伴随的阿贝尔群，记作 $K_G(X)$。若 $G$ 是平凡群，则 $K_G(X)=K(X)$。另一方面，若 $X$ 是单点，则 $K_G(X)\cong R(G)$，即 $G$ 的特征标环。

% [原文 p.45]
设 $E$ 为 $X$ 上的 $G$-向量丛，则 $P(E)$ 是 $G$ 空间。当 $L$ 是 $G$-向量丛时，取 $E=L\oplus1$，则零截面与无穷远截面 $X\to P(E)$ 都是 $G$ 截面。$P(E)$ 上的向量丛 $H$ 也是 $G$ 线丛。由此可见，上一节周期性定理证明中的全部论证都可以在假定 $G$ 作用的情形下进行。因此得到关于 $K_G$ 的周期性定理。

\begin{theorem*}{2.3.1}
  若 $X$ 是 $G$ 空间、$L$ 是 $X$ 上的 $G$ 线丛，则映射 $t\mapsto[H]$ 诱导出 $K_G(X)$ 模的同构映射
  \[
  K_G(X)[t]/((t[L]-1)(t-1))\longrightarrow K_G(P(L\oplus1))
  \]
\end{theorem*}

% [Tutor Note: 省略论证的点拨：定理 2.3.1 没有单列证明，它就在上一段里。只要把 §2.2 的每一步（$P(L\oplus1)$ 的分解、$z$ 与 $H$ 的截面、衔接函数的同伦）逐字换成 $G$ 等变版本即可，所有构造都带 $G$ 作用。]

% [Tutor Note: 记号提醒：原书 p.45 把模写作 $K_G(X)[t]/((t[L]-1)(t-1))$，而 p.42、p.43 写作 $K(X)[t]/((t-1)([L]t-1))$。因 $t$ 与 $[L]$ 可交换，两式生成同一理想，并非笔误，比较时不要误以为差一个符号。]

\section{K 群的一般上同调论性质}
\takeaway{把 $K$ 建成定义在紧空间对上的上同调论：用 $K(X,Y)=\tilde K(X/Y)$ 与约化悬垂 $S$ 定义负次数群 $K^{-n}$，工具是锥 $CX$、悬垂与 smash 积；由此得到对 $(X,Y)$ 的正合列，并最终导出向左无限延伸的长正合列。}

下面，对紧 Hausdorff 空间对 $(X,Y)$ 定义相对 $K$ 群 $K(X,Y)$。这样一来，就可以用完全形式化的方式证明 $K$ 的若干性质。此后给出的定理，其证明是形式化的，因此对满足特定公理的任意“一般上同调论”同样成立。其公理化表述留给读者。

用 $\mathcal{C}$ 表示紧空间的范畴，$\mathcal{C}^+$ 表示带基点的紧空间的范畴，$\mathcal{C}^2$ 表示紧空间对的范畴。函子
\[
\mathcal{C}^2\longrightarrow\mathcal{C}^+
\]
\[
\mathcal{C}\longrightarrow\mathcal{C}^2
\]
定义如下。前者把对 $(X,Y)$ 变为以 $Y/Y$ 为基点的 $X/Y$（若 $Y=\emptyset$，则把 $X/Y$ 视为 $X$ 与单点的不交并）。后者把空间 $X$ 变为对 $(X,\emptyset)$。合成 $\mathcal{C}\to\mathcal{C}^+$ 由 $X\mapsto X^+$ 给出，其中 $X^+$ 表示 $X/\emptyset$。

当 $X$ 属于 $\mathcal{C}^+$ 时，把 $\tilde K(X)$ 定义为映射 $i^*\colon K(X)\to K(x_0)$ 的核。其中 $i\colon\{x_0\}\to X$ 是基点的包含映射。若 $c\colon X\to\{x_0\}$ 是把 $X$ 收缩到基点的映射，
% [原文 p.46]
则 $c^*$ 诱导出分解 $K(X)\cong\tilde K(X)\oplus K(x_0)$。这个分解对 $\mathcal{C}^+$ 的映射显然是自然的。这样一来，$\tilde K$ 是 $\mathcal{C}^+$ 上的函子。又显然有 $K(X)\cong\tilde K(X^+)$。用
\[
K(X,Y)=\tilde K(X/Y)
\]
定义 $K(X,Y)$。特别地，$K(X,\emptyset)\cong K(X)$。由于 $\tilde K$ 是 $\mathcal{C}^+$ 上的函子，可知 $K(X,Y)$ 是 $\mathcal{C}^2$ 中 $(X,Y)$ 的反变函子。

% [Tutor Note: 易错预警：$K(X,Y)$ 的定义是 $\tilde K(X/Y)$，而不是 $K(X)$ 对 $K(Y)$ 的商。只有 $Y\to X$ 是好对（例如 $Y$ 是 $X$ 的形变收缩）时，二者才一致；一般情形下用 $K(X/Y)$ 才是对的。]

接下来，在 $\mathcal{C}^+$ 中引入称为“smash 积”的运算。设 $X,Y\in\mathcal{C}^+$，令
\[
X\wedge Y=(X\times Y)/(X\vee Y)
\]
其中 $X\vee Y=X\times\{y_0\}\cup\{x_0\}\times Y$，$x_0$、$y_0$ 分别是 $X$、$Y$ 的基点。对空间 $X,Y,Z\in\mathcal{C}^+$，有自然同胚
\[
X\wedge(Y\wedge Z)\cong(X\wedge Y)\wedge Z
\]
并且借助这个同胚把这些空间视为同一。

用 $I$ 表示单位区间 $[0,1]$，$\partial I=\{0\}\cup\{1\}$ 表示其边界。把 $I/\partial I\in\mathcal{C}^+$ 作为圆周 $S^1$ 的标准模型。同样，用 $I^n$ 表示 $\R^n$ 中的单位立方体，把 $I^n/\partial I^n$ 作为 $n$ 维球面 $S^n$ 的模型。此时有自然同胚
\[
S^n\cong\underbrace{S^1\wedge S^1\wedge\cdots\wedge S^1}_{n\ \text{个}}
\]
对 $X\in\mathcal{C}^+$，空间 $S^1\wedge X\in\mathcal{C}^+$ 称为 $X$ 的约化悬垂（reduced suspension），多记作 $SX$。$n$ 重悬垂 $SS\cdots SX$（重复 $n$ 次）与 $S^n\wedge X$ 自然同胚，简记作 $S^nX$。

\begin{definition}{2.4.1}
  对 $n\ge0$，定义
  \[
  \tilde K^{-n}(X)=\tilde K(S^nX)\qquad(X\in\mathcal{C}^+)
  \]
  \[
  K^{-n}(X,Y)=\tilde K^{-n}(X/Y)=\tilde K(S^n(X/Y))\qquad((X,Y)\in\mathcal{C}^2)
  \]
  \[
  K^{-n}(X)=K^{-n}(X,\emptyset)=\tilde K(S^n(X^+))\qquad(X\in\mathcal{C})
  \]
\end{definition}

% [Tutor Note: 次数约定：定义 2.4.1 只给出负次数 $K^{-n}$（$n\ge0$），正次数要靠 Bott 周期性 $\tilde K(S^2X)\cong\tilde K(X)$ 反向延拓，即令 $K^{n}=K^{n-2}$。这是后文 $K^*(X)$ 记号的基础，不要以为漏了正次数。]

显然，这些都是各自范畴上的反变函子。

% [原文 p.47]
在继续之前，定义 $X$ 上的锥（cone）为
\[
CX=(I\times X)/(\{0\}\times X)
\]
因此 $C$ 是函子 $C\colon\mathcal{C}\to\mathcal{C}^+$。把 $X$ 与 $CX$ 的子空间 $\{1\}\times X$ 等同。空间 $CX/X=I\times X/\partial I\times X$ 称为 $X$ 的（非约化）悬垂（unreduced suspension）。悬垂是函子 $\mathcal{C}\to\mathcal{C}^+$，而约化悬垂 $S$ 是函子 $\mathcal{C}^+\to\mathcal{C}^+$，此点需要注意。当 $X\in\mathcal{C}^+$ 以 $x_0$ 为基点时，有自然的包含映射
\[
I\cong C\{x_0\}/\{x_0\}\longrightarrow CX/X
\]
在 $CX/X$ 中把 $I$ 收缩成一点所得的商空间恰好是 $SX$。因此由引理 1.4.7，$p\colon CX/X\to SX$ 诱导出同构映射 $K(SX)\cong K(CX/X)$，结果也诱导出同构映射 $\tilde K(SX)\cong K(CX,X)$。这样一来，对约化悬垂与非约化悬垂使用同一记号 $SX$ 也不会出现问题。

% [Tutor Note: 点拨：$CX/X$ 与 $SX$ 只差把子空间 $I$ 收缩成一点，故用引理 1.4.7（紧空间上的同伦等价给出 $K$ 群的同构）即得 $K(SX)\cong K(CX/X)$。这一步是把相对 $K$ 群与约化悬垂接起来的关键，后面的长正合列全靠它。]

当 $(X,Y)\in\mathcal{C}^2$ 时，把 $X$ 与 $CY$ 各自的部分空间 $Y\subset X$ 与 $\{1\}\times Y\subset CY$ 等同，所得的空间记为 $X\cup CY$。以 $CY$ 的基点为 $X\cup CY$ 的基点，则
\[
X\cup CY\in\mathcal{C}^+
\]
$X$ 是 $X\cup CY$ 的子空间，并且有自然同胚
\[
X\cup CY/X\cong CY/Y
\]
此点需要注意。因此，若 $Y\in\mathcal{C}^+$，则
% [原文 p.48]
\[
K(X\cup CY,X)\cong K(CY,Y)\cong\tilde K(SY)=\tilde K^{-1}(Y)
\]
成立。

下面从这个简单的引理开始。

\begin{lemma*}{2.4.2}
  对 $(X,Y)\in\mathcal{C}^2$，有正合列
  \[
  K(X,Y)\xrightarrow{j^*}K(X)\xrightarrow{i^*}K(Y)
  \]
  其中 $i\colon Y\to X$ 与 $j\colon(X,\emptyset)\to(X,Y)$ 是包含映射。
\end{lemma*}

\begin{proof}
  合成 $i^*j^*$ 由合成 $ji\colon(Y,\emptyset)\to(X,Y)$ 诱导，因而通过平凡群 $K(Y,Y)$ 分解。所以 $i^*j^*=0$。其次取 $\xi\in\Ker i^*$。$\xi$ 可以写成 $[E]-[n]$ 的形式，其中 $E$ 是 $X$ 上的向量丛。由 $i^*\xi=0$ 推知，在 $K(Y)$ 中 $[E|_Y]=[n]$。于是对某个整数 $m$ 有
  \[
  (E\oplus m)|_Y=n\oplus m
  \]
  也就是说，得到了 $(E\oplus m)|_Y$ 的一个平凡化 $\alpha$。它确定 $X/Y$ 上的向量丛 $(E\oplus m)/\alpha$，以及元素
  \[
  \eta=[(E\oplus m)/\alpha]-[n\oplus m]\in\tilde K(X/Y)=K(X,Y)
  \]
  此时
  \[
  j^*(\eta)=[E\oplus m]-[n\oplus m]=[E]-[n]=\xi
  \]
  成立。因此 $\Ker i^*=\Image j^*$，正合性得证。
\end{proof}

% [Tutor Note: 点拨：证明中 $i^*j^*=0$ 的理由是，$j^*$ 之后复合的 $ji$ 通过 $(Y,Y)$ 分解，而 $K(Y,Y)=\tilde K(Y/Y)=0$。这是“平凡群”三个字的实际含义。]

\begin{corollary*}{2.4.3}
  若 $(X,Y)\in\mathcal{C}^2$ 且 $Y\in\mathcal{C}^+$（于是把 $Y$ 的基点取作 $X$ 的基点时也有 $X\in\mathcal{C}^+$），则
  \[
  K(X,Y)\longrightarrow\tilde K(X)\longrightarrow\tilde K(Y)
  \]
  是正合列。
\end{corollary*}
% 块尾跨页句：系 2.4.3 的陈述（其正合列与“是正合列”一句在原书 p.49 顶部）由本块写完；ch02c 自 p.49 的“证明”接续，不重述该陈述。

% CHUNK ch02b DONE
